\chapter{Non-Euclidean geometry}\label{geometry:ch:noneuclid}

\begin{center}\itshape
I have created a new, another world out of nothing.
\upshape(J\'anos Bolyai, letter to his father, 3 November 1823)
\end{center}

Draw a line and a point outside it. How many lines through that point never
meet the first line? In the Euclidean plane the answer is one. On a sphere,
where lines are great circles, the answer is none. In the hyperbolic plane
there are infinitely many. The choice changes distances, angle sums and areas.

Euclid's fifth postulate describes when two lines, cut by a third, must meet.
In the usual formulation by parallels, it says that exactly one parallel
passes through a point outside a line. Attempts to derive it from the other
postulates continued for centuries. Saccheri's 1733 book pursued consequences
of denying it, hoping to reach a contradiction. The consequences were
unfamiliar, but unfamiliarity was not a contradiction.

Lobachevsky published his hyperbolic geometry in 1829. Bolyai developed it
independently and published it in 1831--32, as an appendix to his father's
textbook. His letter of 1823 supplies the epigraph. Gauss had also considered
the subject privately. Beltrami's models of 1868 and Poincar\'e's work on
Fuchsian groups in 1882 made the geometry accessible through familiar
coordinates. A model turns a proposed geometry into explicit calculations.

For the sphere, a triangle's area is its angle excess over $\pi$.
For the hyperbolic plane, the area is its angle deficit below $\pi$.
Chapter~9 has already proved the hyperbolic formula, distance and isometry
results. Here we prove the spherical formula and compare the two geometries
through calculations. The prerequisites are arc length, integration and
complex numbers, together with those results from Chapter~9.

\begin{think}
The North Pole and two points on the equator, separated by a quarter-turn,
bound a triangle made of three great-circle arcs. What are its three angles?
What fraction of the sphere does it occupy?
\end{think}

\section{Spherical geometry}

A great circle is the intersection of the unit sphere $S^2\subset\mathbb R^3$
with a plane through $0$. Distinct great circles meet in two antipodal points,
because their two planes meet in a line through $0$. Thus they supply no
parallels. The spherical distance between $u,v\in S^2$ is
$\arccos(u\cdot v)$, the length of their shorter great-circle arc.
For antipodal points the distance is $\pi$, but the shorter arc is not unique.
We do not need a general shortest-path theorem for the calculations below.

A spherical triangle here means a nondegenerate convex triangle contained
in an open hemisphere, bounded by the shorter great-circle arcs between
its vertices. Its sides have lengths in $(0,\pi)$ and its interior angles
lie in $(0,\pi)$. These choices matter: the same three great circles also
bound other regions, to which the displayed triangle formula would not
apply with the same angles.

\begin{theorem}[Girard, 1629]\label{geometry:thm:girard}
On the unit sphere, a spherical triangle with interior angles
$\alpha,\beta,\gamma$ has area $\alpha+\beta+\gamma-\pi$.
\end{theorem}
\begin{proof}
A lune is the region between two great semicircles with the same antipodal
endpoints. If its angle is $\theta$, its area is $2\theta$.
For example, spherical coordinates with those endpoints as poles give
\[
 \operatorname{area}(\text{lune})
 =\int_0^\theta\int_0^\pi\sin t\,\dd t\,\dd\phi=2\theta.
\]

Let $T$ be our triangle and $-T$ its antipodal image. At each vertex, take
the lune of the corresponding interior angle that contains $T$, and also
its antipodal lune containing $-T$. There are six lunes, of total area
$4(\alpha+\beta+\gamma)$. The three great circles cut the sphere into
$T$, $-T$ and six further triangles. The interiors of $T$ and $-T$ belong
to three of the chosen lunes each; each other triangular region belongs
to one. Boundaries have area zero. Since the sphere has area $4\pi$,
\[
 4(\alpha+\beta+\gamma)
 =4\pi+2\operatorname{area}(T)+2\operatorname{area}(-T)
 =4\pi+4\operatorname{area}(T).
\]
Rearranging proves the formula.
\end{proof}

\begin{figure}[ht]
\centering
\fig{geometry-ch10-sphere}
\caption{The spherical triangle with vertices $e_1,e_2,e_3$ has three right angles. Its area is $\frac{3\pi}2-\pi=\frac\pi2$, one eighth of the sphere.}
\label{geometry:fig:sphere}
\end{figure}

The coordinate octant has three right angles and area $\pi/2$, as the
formula predicts. Its excess is not an error in measurement; the excess
measures the area. For a sphere of radius $R$, lengths scale by $R$ and
areas by $R^2$, so the area is $R^2(\alpha+\beta+\gamma-\pi)$.

Distances also obey a different cosine law. The dot product gives it
directly, without a new trigonometric identity to memorise.

\begin{proposition}[spherical law of cosines]\label{geometry:prop:sphcos}
Let $a,b,c$ be the side lengths opposite the vertices $A,B,C$ of a
spherical triangle on the unit sphere, and let $\gamma$ be its angle at $C$.
Then
\[
 \cos c=\cos a\cos b+\sin a\sin b\cos\gamma.
\]
\end{proposition}
\begin{proof}
Regard $A,B,C$ as unit vectors. Their dot products are
$B\cdot C=\cos a$, $A\cdot C=\cos b$ and $A\cdot B=\cos c$.
Project $A$ and $B$ onto the tangent plane at $C$, and normalise:
\[
 t_A=\frac{A-(A\cdot C)C}{\sin b},\qquad
 t_B=\frac{B-(B\cdot C)C}{\sin a}.
\]
The denominators are positive because $a,b\in(0,\pi)$. These unit tangent
vectors point along the sides from $C$ towards $A$ and $B$, so their dot
product is $\cos\gamma$. Expanding it gives
\[
 \cos\gamma
 =\frac{A\cdot B-(A\cdot C)(B\cdot C)}{\sin a\sin b}
 =\frac{\cos c-\cos a\cos b}{\sin a\sin b}.
\]
Multiplication by the denominator gives the stated law.
\end{proof}

For a right triangle it becomes $\cos c=\cos a\cos b$.
Expanding $\cos x=1-x^2/2+O(x^4)$ recovers
$c^2=a^2+b^2+O(a^4+b^4)$ for small sides. The Euclidean law is the local
approximation; the higher-order terms detect curvature.

\section{The hyperbolic plane}

In the upper half-plane, the length element is $|\dd z|/y$.
Moving higher makes a given Euclidean displacement shorter in this metric.
The h-lines are vertical lines and semicircles perpendicular to the real
axis. Chapter~9 proves that their arcs are shortest paths and that
\[
 \cosh d(z,w)=1+\frac{|z-w|^2}
 {2\operatorname{Im}z\operatorname{Im}w}
\]
(Theorem~\ref{geometry:thm:hypdist}). It also proves
\[
 \operatorname{area}(T)=\pi-\alpha-\beta-\gamma
\]
for hyperbolic triangles, allowing ideal vertices with angle $0$
(Theorem~\ref{geometry:thm:hyparea}). A nondegenerate triangle has positive
area. Its angle sum is therefore less than $\pi$. Its area is at most
$\pi$, with equality precisely when all three vertices are ideal.

The Cayley map transports the metric to the unit disc as
$2|\dd w|/(1-|w|^2)$ (Proposition~\ref{geometry:prop:disk}).
The boundary circle is infinitely far away. A point near it can look
Euclidean-close to another point while remaining a large hyperbolic
distance away. In polar distance coordinates, the circumference of a
circle of radius $r$ is $2\pi\sinh r$. Compare this with $2\pi\sin r$
on the unit sphere and $2\pi r$ in the plane; the three agree to first
order at $r=0$ and then separate.

The disc metric is unchanged by rotations and by complex conjugation.
Reflection in a circular h-line is inversion in its circle; reflection
in a diameter is ordinary reflection in that diameter.
Repeated reflections in a hyperbolic triangle of angles
$\pi/2,\pi/3,\pi/7$ produce the tiling in
Figure~\ref{geometry:fig:tiling}. The existence of that reflection tiling
is a quoted fact about triangle groups. Four triangles fit around the
right-angle vertex, six around the $\pi/3$ vertex and fourteen around
the $\pi/7$ vertex. Each triangle has area
\[
 \pi-\frac\pi2-\frac\pi3-\frac\pi7=\frac\pi{42}.
\]

\begin{figure}[ht]
\centering
\fig{geometry-ch10-tiling}
\caption{The Poincar\'e disk tiled by hyperbolic triangles with angles $\frac\pi2$, $\frac\pi3$, $\frac\pi7$, alternately coloured. All the triangles are
congruent: they only look smaller near the boundary, which is infinitely far away.}
\label{geometry:fig:tiling}
\end{figure}

Among hyperbolic triangles with angles $\pi/p,\pi/q,\pi/r$ for integers
$p,q,r\ge2$, this area is smallest. Exercise~\ref{geometry:ex:tilings}
proves the numerical bound from the angle formula. The reflection-group
theorem supplying the tilings is separate from that calculation.

Coxeter sent Escher a paper with a hyperbolic tiling in 1958. Escher's
reply described the effect as ``quite a shock''. His subsequent
\emph{Circle Limit} woodcuts use the same way of fitting infinitely many
congruent pieces inside a Euclidean circle. The pieces shrink in the
picture because the metric scale grows near the boundary.

\begin{remark}
A quoted theorem of Hurwitz bounds the number of holomorphic automorphisms
of a compact Riemann surface of genus $g\ge2$ by $84(g-1)$.
The constant comes from the orientation-preserving triangle group of type
$(2,3,7)$. Its fundamental region has twice the triangle's area,
$\pi/21$, and $4\pi(g-1)/(\pi/21)=84(g-1)$.
This is an application of triangle groups, not a consequence of the area
calculation alone.
\end{remark}

\begin{figure}[ht]
\centering
\fig{v9-app-i-polar-circles}
\caption{In the Poincar\'e disc, circles at hyperbolic distances $r=0.5,1,2,3$ from $0$ have Euclidean radii $\tanh(r/2)$. Their Euclidean spacing shrinks towards the boundary while their hyperbolic circumference $2\pi\sinh r$ grows. The dashed boundary is not part of the plane.}
\label{geometry:fig:v9-polar-circles}
\end{figure}

For subdivisions, it can be useful to draw h-lines straight. The radial change
of coordinates $k(w)=2w/(1+|w|^2)$ sends the Poincar\'e disc bijectively
to the \emph{Klein disc}. Indeed, its radial coordinate $2\rho/(1+\rho^2)$
increases from $0$ to $1$. A circular h-line has equation
$|w|^2-2a\cdot w+1=0$ for a real plane vector $a$; this becomes
$a\cdot k(w)=1$, a straight line. Diameters stay diameters.
Thus h-line polygons become ordinary convex polygons in these coordinates,
which justifies the fan subdivision in the next figure. This drawing
preserves incidence, but it does not preserve angles.

\begin{figure}[ht]
\centering
\fig{v9-app-i-triangle-area}
\caption{A fan from one vertex divides a convex hyperbolic pentagon into three triangles. Summing their areas gives $3\pi$ minus the sum of the pentagon's interior angles. The drawn pentagon is in the Klein disc, where h-lines appear straight; its Euclidean angles do not represent hyperbolic angles.}
\label{geometry:fig:v9-triangle-area}
\end{figure}

The angle formula measures a difference that no change of drawing can erase.
A small spherical triangle has a small positive excess; a small hyperbolic
triangle has a small positive deficit. Their local Euclidean appearance does
not make the two geometries identical. Appendix~J turns to a different
measurement question: how much area can a boundary of fixed length enclose?

\section{Exercises}

\begin{exercise}[$\star$]\label{geometry:ex:lune}
Find the areas of a spherical triangle with three right angles and of one
with three angles $2\pi/3$. In each case, exhibit a tiling of the unit sphere
by congruent copies and count them.
\end{exercise}
\begin{exercise}[$\star$]\label{geometry:ex:sphpythag}
Obtain the right-triangle formula $\cos c=\cos a\cos b$ from the spherical
cosine law. Prove $c^2=a^2+b^2+O(a^4+b^4)$ as $a,b\to0$.
\end{exercise}
\begin{exercise}[$\star$]\label{geometry:ex:sphcircle}
On the unit sphere, compute the length of the circle at distance
$r\in(0,\pi)$ from the North Pole. Show it is less than $2\pi r$.
\end{exercise}
\begin{exercise}[$\star$]\label{geometry:ex:hpdist}
Determine the h-line through $1+i$ and $3+i$, and find the distance between
these points. Compare it with the length $2$ of the horizontal segment.
\end{exercise}
\begin{exercise}[$\star$]\label{geometry:ex:inversioniso}
Prove that $z\mapsto1/\overline z$ is an isometry of $\mathbb H$ fixing
$|z|=1$ pointwise. Deduce the same facts for inversion in any semicircle
centred on the real axis.
\end{exercise}
\begin{exercise}[$\star\star$]\label{geometry:ex:circles}
Find the Euclidean radius of a hyperbolic circle of radius $r>0$ about $0$
in the disc model. Compute its hyperbolic length and the area it encloses.
\end{exercise}
\begin{exercise}[$\star\star$]\label{geometry:ex:hypcos}
For a hyperbolic triangle with side lengths $a,b,c$ and angle $\gamma$
opposite $c$, prove
$\cosh c=\cosh a\cosh b-\sinh a\sinh b\cos\gamma$.
State its right-triangle case.
\end{exercise}
\begin{exercise}[$\star$]\label{geometry:ex:ideal}
Show that any two triangles with three distinct ideal vertices are
congruent under a hyperbolic isometry, allowing orientation reversal.
\end{exercise}
\begin{exercise}[$\star\star$]\label{geometry:ex:parallelangle}
Let $F=i$ and $P=e^{i\theta}$ with $0<\theta<\pi/2$.
The unit semicircle through $P$ meets the imaginary axis perpendicularly
at $F$. Show that the h-line from $P$ towards the ideal endpoint $0$
makes angle $\theta$ with the arc $PF$ at $P$.
If $d=d(P,F)$, prove $\tan(\theta/2)=e^{-d}$.
\end{exercise}
\begin{exercise}[$\star$]\label{geometry:ex:tilings}
Suppose a hyperbolic triangle has angles $\pi/p,\pi/q,\pi/r$, where
$p,q,r\ge2$ are integers. Prove $1/p+1/q+1/r<1$ and show its area is
at least $\pi/42$. Determine every equality case.
\end{exercise}
\begin{exercise}[$\star\star$]\label{geometry:ex:quadrilateral}
Prove that a convex hyperbolic quadrilateral has angle sum less than $2\pi$.
Why does this rule out a rectangle with four right angles?
\end{exercise}
\begin{exercise}[$\star\star\star$]\label{geometry:ex:horocycle}
Show that, for every fixed $z=x+i$ on the line $y=1$,
$d(z,iT)-\log T\to0$ as $T\to\infty$.
This describes a horocycle centred at the ideal point $\infty$.
Compute both its arc length from $i$ to $x+i$ and the hyperbolic distance
between those endpoints. How different can the two lengths become?
\end{exercise}

\begin{exercise}[$\star\star$]\label{geometry:ex:v9-caps-and-balls}
\emph{Area from circles.} Compute the area of a spherical cap of radius
$r\in(0,\pi)$ on the unit sphere and rederive the area of a hyperbolic
ball of radius $r>0$ by integrating its circumference.
Compare both with $\pi r^2$ for small $r$, through the term in $r^4$.
For $r<\pi$, prove the strict ordering of these three areas.
\end{exercise}
\begin{exercise}[$\star\star$]\label{geometry:ex:v9-offset-circle}
\emph{Where a hyperbolic centre appears.} In $\mathbb H$, show that the
hyperbolic circle of radius $r>0$ about $x_0+iy_0$ is a Euclidean circle
with centre $x_0+iy_0\cosh r$ and radius $y_0\sinh r$.
Find its highest and lowest points and their hyperbolic distance apart.
Explain why its Euclidean and hyperbolic centres differ.
\end{exercise}
\begin{exercise}[$\star\star$]\label{geometry:ex:v9-hyperbolic-polygons}
\emph{Angle deficit for a polygon.} Let a convex hyperbolic $n$-gon,
$n\ge3$, have finite vertices and interior angles $\alpha_1,\ldots,\alpha_n$.
Prove its area is $(n-2)\pi-\sum_j\alpha_j$.
Deduce the necessary condition on $n$ for all angles to be right angles,
and find the area of a right-angled hexagon.
State and justify the area of an ideal convex $n$-gon.
\end{exercise}
\begin{exercise}[$\star\star$]\label{geometry:ex:v9-rigid-equilateral}
\emph{Angles fix a scale.} An equilateral hyperbolic triangle has side
length $s>0$ and angle $\alpha$. Prove
$\cos\alpha=\cosh s/(1+\cosh s)$.
Show that two such triangles with the same angles have the same side
length. Determine the limits of $\alpha$ and the area as $s\to0$ and
as $s\to\infty$. Why can there be no map of the hyperbolic plane that
preserves angles and multiplies every distance by a fixed $c\ne1$?
\end{exercise}

\section{Hints}

\begin{hint}{geometry:ex:lune}
Use the coordinate octants for the first tiling and radial projection of
the faces of a regular tetrahedron for the second.
\end{hint}
\begin{hint}{geometry:ex:sphpythag}
First bound $c\le a+b$. Then the remainder $O(c^4)$ is controlled by
$O(a^4+b^4)$.
\end{hint}
\begin{hint}{geometry:ex:sphcircle}
All points of the circle have height $\cos r$.
\end{hint}
\begin{hint}{geometry:ex:hpdist}
The semicircle's centre is real and equidistant from the two points.
\end{hint}
\begin{hint}{geometry:ex:inversioniso}
Compose $z\mapsto-1/z$ with $z\mapsto-\overline z$, then conjugate
by $z\mapsto c+rz$ for $c\in\mathbb R$, $r>0$.
\end{hint}
\begin{hint}{geometry:ex:circles}
Use $d(0,w)=2\operatorname{artanh}|w|$ and integrate the disc metric
and its area element in polar coordinates.
\end{hint}
\begin{hint}{geometry:ex:hypcos}
Place the vertex of angle $\gamma$ at $0$ in the disc. The other two
vertices have moduli $\tanh(b/2)$ and $\tanh(a/2)$.
\end{hint}
\begin{hint}{geometry:ex:ideal}
A real fractional linear map moves three boundary points to $0,1,\infty$.
Check its determinant's sign to decide whether conjugation is needed.
\end{hint}
\begin{hint}{geometry:ex:parallelangle}
The circle through $P$ and $0$ has real centre $c$ with $|P-c|=c$.
Use the isosceles triangle $0,c,P$ and then compute $\cosh d$.
\end{hint}
\begin{hint}{geometry:ex:tilings}
Order $p\le q\le r$. Separate $p\ge3$, $p=2,q\ge4$, and $p=2,q=3$.
\end{hint}
\begin{hint}{geometry:ex:quadrilateral}
A diagonal gives two triangles. Add their angles.
\end{hint}
\begin{hint}{geometry:ex:horocycle}
Compute $\cosh d(x+i,iT)$ exactly before taking the limit.
\end{hint}
\begin{hint}{geometry:ex:v9-caps-and-balls}
Integrate $2\pi\sin t$ or $2\pi\sinh t$ from $0$ to $r$.
Use $\sin t<t<\sinh t$ for the comparison.
\end{hint}
\begin{hint}{geometry:ex:v9-offset-circle}
Set the distance formula equal to $\cosh r$ and complete a square in $y$.
\end{hint}
\begin{hint}{geometry:ex:v9-hyperbolic-polygons}
Draw all diagonals from one vertex, as in
Figure~\ref{geometry:fig:v9-triangle-area}. For ideal vertices use ideal
triangles directly; their area formula was proved in Chapter~9.
\end{hint}
\begin{hint}{geometry:ex:v9-rigid-equilateral}
Put $a=b=c=s$ in the hyperbolic cosine law, then use
$\sinh^2s=(\cosh s-1)(\cosh s+1)$.
\end{hint}

\section{Solutions}

\begin{solution}{geometry:ex:lune}
\idea The angle excess gives an area; a symmetric solid shows whether those
triangles really tile.
\writeup Three right angles give area $3\pi/2-\pi=\pi/2$.
The coordinate planes cut the sphere into eight congruent such octants.
Three angles $2\pi/3$ give area $2\pi-\pi=\pi$.
Project the four faces of a regular tetrahedron radially from its centre
onto the sphere. The edges become shorter great-circle arcs. At each
vertex, symmetry makes the three angles equal, so each is $2\pi/3$.
The four spherical faces cover the sphere and have the stated area.
\end{solution}
\begin{solution}{geometry:ex:sphpythag}
\idea Keep the first nonconstant term of the cosine law and control the
error before solving for $c^2$.
\writeup For $\gamma=\pi/2$, the cosine term vanishes and
$\cos c=\cos a\cos b$. Expanding both sides gives
\[
 1-c^2/2+O(c^4)=1-(a^2+b^2)/2+O(a^4+b^4),
\]
since $a^2b^2\le(a^4+b^4)/2$. For small $a,b$ we have $a+b<\pi$, and
$\cos c=\cos a\cos b\ge\cos(a+b)$ because $\sin a\sin b\ge0$.
Cosine decreases on $[0,\pi]$, so $c\le a+b$.
Thus $c^4\le8(a^4+b^4)$, and rearrangement gives the claimed estimate.
\end{solution}
\begin{solution}{geometry:ex:sphcircle}
\idea A distance circle about the pole is a latitude circle in space.
\writeup For $N=(0,0,1)$, the condition $d(u,N)=r$ says
$u\cdot N=\cos r$. The circle has height $\cos r$ and Euclidean radius
$\sqrt{1-\cos^2r}=\sin r$. Its arc length is therefore $2\pi\sin r$.
Since $\sin r<r$ for $r>0$, this is less than $2\pi r$.
\end{solution}
\begin{solution}{geometry:ex:hpdist}
\idea The symmetry of the two points locates the circle's centre.
\writeup Their perpendicular bisector is $x=2$, so the required real
centre is $2$ and the radius is $\sqrt2$. The h-line is
$(x-2)^2+y^2=2$, $y>0$.
The distance formula gives $\cosh d=1+4/2=3$, hence
$d=\log(3+2\sqrt2)\approx1.763<2$.
The horizontal segment has length $\int_1^3\dd x=2$.
The h-line arc rises above it, where the length scale $1/y$ is smaller.
\end{solution}
\begin{solution}{geometry:ex:inversioniso}
\idea Known isometries remain isometries when composed or conjugated.
\writeup The matrix $\left(\begin{smallmatrix}0&-1\\1&0\end{smallmatrix}\right)$
has determinant $1$, so $z\mapsto-1/z$ is an isometry by Chapter~9.
Reflection $z\mapsto-\overline z$ is also an isometry. Their composition
is $z\mapsto1/\overline z$, fixing every $z$ with $|z|=1$.
For centre $c\in\mathbb R$ and radius $r>0$, conjugate by the isometry
$h(z)=c+rz$. The result is
$h(1/\overline{h^{-1}(z)})=c+r^2/\overline{z-c}$,
the inversion in the required semicircle. Conjugation also carries its
pointwise fixed set to that semicircle.
\end{solution}
\begin{solution}{geometry:ex:circles}
\idea Rotational symmetry lets one distance calculation determine the
whole circle.
\writeup The disc distance formula gives $d(0,w)=2\operatorname{artanh}|w|$.
Thus the Euclidean radius is $\rho=\tanh(r/2)$.
On $|w|=\rho$ the metric factor is constant, so its length is
\[
 \frac{2}{1-\rho^2}\,2\pi\rho=2\pi\sinh r.
\]
The area element is $4\,\dd A/(1-|w|^2)^2$. Consequently its area is
\[
 \int_0^\rho\frac{8\pi t}{(1-t^2)^2}\,\dd t
 =\frac{4\pi\rho^2}{1-\rho^2}=4\pi\sinh^2(r/2).
\]
For small $r$ these approach $2\pi r$ and $\pi r^2$ respectively.
\end{solution}
\begin{solution}{geometry:ex:hypcos}
\idea Put the angle at the disc's centre so that its two sides are radii.
\writeup Use an isometry to set $C=0$, $A=u=\tanh(b/2)$ and
$B=ve^{i\gamma}$ with $v=\tanh(a/2)$. Angles in this model are Euclidean.
The distance formula in Proposition~\ref{geometry:prop:disk} gives
\[
 \cosh c=1+\frac{2(u^2+v^2-2uv\cos\gamma)}{(1-u^2)(1-v^2)}
 =\frac{(1+u^2)(1+v^2)-4uv\cos\gamma}{(1-u^2)(1-v^2)}.
\]
Now $(1+u^2)/(1-u^2)=\cosh b$ and $2u/(1-u^2)=\sinh b$;
the corresponding identities for $v$ give the required cosine law.
For $\gamma=\pi/2$ it reads $\cosh c=\cosh a\cosh b$.
\end{solution}
\begin{solution}{geometry:ex:ideal}
\idea Three distinct boundary points can be normalised by one fractional
linear transformation.
\writeup By Theorem~\ref{geometry:thm:mobius}(2), a fractional linear map
$g(z)=(az+b)/(cz+d)$ sends the three vertices to $0,1,\infty$.
Its coefficients can be chosen real because the vertices are real or
infinite. Write $\Delta=ad-bc\ne0$.
If $\Delta>0$, scaling all coefficients by $1/\sqrt\Delta$ gives an
isometry of $\mathbb H$ with the required boundary values.
If $\Delta<0$, the map $-g$ has positive determinant and sends the
vertices to $0,-1,\infty$; composing with $z\mapsto-\overline z$
sends them to $0,1,\infty$. Isometries carry h-lines to h-lines,
so in either case they carry the entire triangle to the standard one.
Composing the normalisations for two triangles proves congruence.
\end{solution}
\begin{solution}{geometry:ex:parallelangle}
\idea The Euclidean radii of two circles also determine their intersection
angle in the conformal model.
\writeup The circle through $0$ and $P$ has real centre $c$ with
$|P-c|^2=c^2$. Since $|P|=1$, this gives
$c=1/(2\cos\theta)>0$. The triangle $0,c,P$ is isosceles, and
its angles at $0$ and $P$ are both $\theta$.
The radii at $P$ of this circle and the unit circle are $P-c$ and $P$.
The acute angle between them, and hence between the two tangents, is
$\theta$. Choosing the rays towards $0$ and $F$ gives this acute angle.

The distance formula gives
$\cosh d=1+(2-2\sin\theta)/(2\sin\theta)=1/\sin\theta$.
Put $u=\tan(\theta/2)\in(0,1)$. Then
$\cosh d=(u+u^{-1})/2$.
The positive root $e^d\ge1$ of $X+X^{-1}=u+u^{-1}$ is $u^{-1}$,
so $u=e^{-d}$. The angle tends to $\pi/2$ as $d\to0$, and to $0$
as $d\to\infty$.
\end{solution}
\begin{solution}{geometry:ex:tilings}
\idea Maximise the reciprocal sum while keeping it below $1$.
\writeup Positive triangle area gives $1/p+1/q+1/r<1$.
Order $p\le q\le r$. If $p\ge3$, either $q=3$ and $r\ge4$,
giving a sum at most $11/12$, or $q\ge4$, giving a sum at most
$1/3+1/4+1/4=5/6$. Thus the deficit is at least $1/12$.
If $p=2$ and $q\ge4$, then $q=4$ forces $r\ge5$, with sum at most
$19/20$; if $q\ge5$, the sum is at most $1/2+1/5+1/5=9/10$.
The deficit in this case is at least $1/20$.
The remaining possible case is $p=2,q=3$, where $r\ge7$ and the
smallest deficit is $1/42$, attained only at $r=7$.
The case $p=q=2$ is impossible. Multiplying the least deficit by $\pi$
proves the area bound and the equality condition $\{p,q,r\}=\{2,3,7\}$.
\end{solution}
\begin{solution}{geometry:ex:quadrilateral}
\idea The angle sum is additive under a diagonal subdivision.
\writeup Convexity keeps diagonal $AC$ inside $ABCD$.
The two nondegenerate triangles have angle sums less than $\pi$ each.
At $A$ and $C$ their angles add to the corresponding quadrilateral
angles; the other two are unchanged. The total is less than $2\pi$.
Four right angles would total $2\pi$, so such a rectangle cannot exist.
The same addition gives area $2\pi$ minus the quadrilateral's angle sum.
\end{solution}
\begin{solution}{geometry:ex:horocycle}
\idea A distant point up the imaginary axis makes a fixed horizontal
separation negligible.
\writeup For each fixed real $x$, the distance formula gives
\[
 \cosh d(x+i,iT)=\frac T2+\frac{x^2+1}{2T}.
\]
As $u\to\infty$, $\operatorname{arcosh}u=\log(2u)+O(u^{-2})$.
Therefore
\[
 d(x+i,iT)-\log T
 =\log\left(1+\frac{x^2+1}{T^2}\right)+O_x(T^{-2})\to0.
\]
This is a limit for each fixed $x$, not a uniform assertion over the
whole line. Along $y=1$, the arc length is $\int_0^x|\dd t|=|x|$.
For the endpoints, $\cosh d=1+x^2/2$ and hence
$d(i,x+i)=2\operatorname{arsinh}(|x|/2)$.
It grows logarithmically for large $|x|$, while the horocycle arc grows
linearly. For $x=0$, both lengths are zero.
\end{solution}
\begin{solution}{geometry:ex:v9-caps-and-balls}
\idea Add thin rings, using distance from the centre as the radial variable.
\writeup On the sphere, the polar parametrisation has area element
$\sin t\,\dd t\,\dd\theta$. Thus
$A_S(r)=\int_0^r2\pi\sin t\,\dd t=2\pi(1-\cos r)$.
In the hyperbolic disc, substituting $\rho=\tanh(t/2)$ into the polar
area element gives $\sinh t\,\dd t\,\dd\theta$. Thus
$A_H(r)=2\pi(\cosh r-1)=4\pi\sinh^2(r/2)$.
Taylor expansion gives
\[
 A_S(r)=\pi r^2-\frac\pi{12}r^4+O(r^6),\qquad
 A_H(r)=\pi r^2+\frac\pi{12}r^4+O(r^6).
\]
For $0<t<\pi$, $\sin t<t<\sinh t$.
Integrating the strict inequalities over $(0,r)$ proves
$A_S(r)<\pi r^2<A_H(r)$ for every $0<r<\pi$.
\end{solution}
\begin{solution}{geometry:ex:v9-offset-circle}
\idea Convert constant hyperbolic distance into a quadratic equation in
Euclidean coordinates.
\writeup For $z=x+iy$ and $z_0=x_0+iy_0$, the distance equation is
\[
 (x-x_0)^2+(y-y_0)^2=2yy_0(\cosh r-1).
\]
Completing the square gives
\[
 (x-x_0)^2+(y-y_0\cosh r)^2=y_0^2\sinh^2r.
\]
Its lowest height is $y_0(\cosh r-\sinh r)=y_0e^{-r}>0$,
so the entire circle stays in $\mathbb H$.
Its highest point has height $y_0e^r$. Both lie on the h-line $x=x_0$,
and their distance is
$\log((y_0e^r)/(y_0e^{-r}))=2r$.
The hyperbolic centre has height $y_0$, while the Euclidean centre has
height $y_0\cosh r>y_0$. The metric gives a smaller length to the same
Euclidean displacement higher up; the circle therefore extends farther
upward in Euclidean coordinates.
\end{solution}
\begin{solution}{geometry:ex:v9-hyperbolic-polygons}
\idea A fan has $n-2$ triangles and introduces no extra boundary angle.
\writeup Join one vertex to all nonadjacent vertices.
Convexity places these diagonals inside the polygon, dividing it into
$n-2$ triangles with disjoint interiors. Summing their triangle formulas
adds $(n-2)\pi$. At each vertex, the incident triangle angles add to
the polygon's interior angle. Areas add as well, proving
$A=(n-2)\pi-\sum_j\alpha_j$.
If all angles are $\pi/2$, this becomes
$A=(n-4)\pi/2>0$, so necessarily $n\ge5$.
A right-angled hexagon has area $\pi$.

An ideal convex polygon has distinct cyclically ordered boundary
vertices. A fan divides it into $n-2$ ideal triangles. Each has area
$\pi$ by Chapter~9; boundaries have area zero even though their length
is infinite. Hence its area is $(n-2)\pi$.
\end{solution}
\begin{solution}{geometry:ex:v9-rigid-equilateral}
\idea Equal sides reduce the cosine law to a strictly monotone relation
between one length and one angle.
\writeup Put $a=b=c=s$ in Exercise~\ref{geometry:ex:hypcos}.
With $u=\cosh s>1$, the equation becomes
$u=u^2-(u^2-1)\cos\alpha$. Thus
$\cos\alpha=(u^2-u)/(u^2-1)=u/(u+1)$.
The function $u/(u+1)$ is strictly increasing for $u>1$, as is
$\cosh s$ for $s>0$. Equal angles therefore imply equal $s$.
As $s\to0$, $u\to1$, so $\alpha\to\pi/3$ and
$A=\pi-3\alpha\to0$. As $s\to\infty$, $u\to\infty$, so
$\alpha\to0$ and $A\to\pi$.

Suppose a map preserves angles and multiplies every distance by
$c>0$. It takes a geodesic segment to the geodesic segment between the
images: equality in the distance triangle inequality along the segment
is preserved, and shortest arcs are unique. Thus an equilateral triangle
of side $s$ maps to one of side $cs$ with the same angles.
The relation just proved forces $cs=s$, hence $c=1$.
A nonpositive factor cannot define such an angle-preserving map with
distinct image points, so no omitted scale provides an exception.
\end{solution}

\section*{Sources and further reading}

\begin{itemize}
\item The MacTutor biography of J\'anos Bolyai (University of St Andrews): the letter of 3 November 1823 (epigraph), the Appendix sent to Gauss in 1831, Gauss's
reply ``to praise it would amount to praising myself'', and Bolyai's discovery in 1848 that Lobachevsky had published similar work in 1829.
\item G.~Saccheri, \emph{Euclides ab omni naevo vindicatus}, Milan, 1733; E.~Beltrami, \emph{Saggio di interpretazione della geometria non-euclidea}, 1868;
H.~Poincar\'e, \emph{Th\'eorie des groupes fuchsiens}, Acta Mathematica 1 (1882).
\item A.~Girard, \emph{Invention nouvelle en l'alg\`ebre}, Amsterdam, 1629.
\item D.~Dunham, \emph{Hyperbolic art and the poster design}, Mathematics Awareness Month 2003 essay (Coxeter's 1957 paper \emph{Crystal symmetry and its
generalizations}, sent to Escher in 1958, and Escher's ``quite a shock'').
\item J.~W.~Anderson, \emph{Hyperbolic Geometry}, Springer, 2nd ed., 2005; A.~F.~Beardon, \emph{The Geometry of Discrete Groups}, Springer, 1983 (the half-plane and
disk models, triangle groups).
\item J.~Stillwell, \emph{Geometry of Surfaces}, Springer, 1992; M.~J.~Greenberg, \emph{Euclidean and Non-Euclidean Geometries}, Freeman, 4th ed., 2008 (the history
of the parallel postulate).
\end{itemize}
